\section{Higher-dimensional extension: available energy}
\label{sec:3d}

The enabling property of unbinned OmniFold is that adding an observable is
adding an input feature: the step-1/step-2 classifiers simply gain a column,
with no migration-matrix or binning penalty. There is no practical D'Agostini
IBU analog, which makes this the primary physics motivation for the technique.
We exercise it by extending the measurement to
$d^{3}\sigma/(d\pT\,d\pPar\,d\Eavail)$ with available energy as a third axis.

The significance is clear from the measurement record: every published MINERvA
differential cross section has unfolded at most two kinematic variables
simultaneously (Fig.~\ref{fig:landscape}, Table~\ref{tab:landscape}), the
double-differential measurements all relying on binned IBU. A simultaneous
\emph{three}-observable unfold of MINERvA data is, to our knowledge,
without precedent, and is a direct consequence of the unbinned formulation.

\begin{figure}[H]
  \centering
  \includegraphics[width=0.97\textwidth]{minerva_unfolding_landscape}
  \caption{Dimensionality of the full set of MINERvA differential cross-section
  measurements over time (number of kinematic variables unfolded
  \emph{simultaneously}; a double-differential $d^{2}\sigma/dx\,dy$ counts as
  two), coloured by channel. All $N=30$ prior measurements use binned D'Agostini
  IBU and reach at most two variables; this work's unbinned OmniFold adds a
  third. Source: the MINERvA data-release page; full list in
  Table~\ref{tab:landscape}. Dimensionality is classified from each analysis's
  headline differential result; flux, reconstruction, generator-tuning and
  total-only entries are excluded.}
  \label{fig:landscape}
\end{figure}



\subsection{Definitions and pipeline}
\label{sec:3d-method}
Available energy follows the MINERvA low-recoil convention
(Ref.~\cite{MINERvA:2023avz}, Eq.~4): truth
$\Eavail=\sum_{p}\mathrm{KE}+\sum_{\pi^{\pm}}\mathrm{KE}+\sum_{\pi^{0},\gamma}E$
(excluding neutrons and the muon), via \texttt{CVUniverse::GetEAvailableTrue()};
reco $\Eavail$ is the calorimetric tracker$+$ECAL estimate scaled by $1.17$,
\texttt{CVUniverse::NewEavail()}. The event loop evaluates both accessors
over all twelve ME~FHC playlists, carrying the $\Eavail$
branches alongside the $\pT$/$\pPar$ branches. The unfolding driver is
the 2D pipeline with a three-column feature stack
$(\pT,\pPar,\Eavail)$; the cross-section extraction generalises to a 3D Jacobian
$\Delta\pT\,\Delta\pPar\,\Delta\Eavail$ with the flux $\Phi$ still applied per
$\pT$ (broadcast over $\pPar$ and $\Eavail$). The $\Eavail$ binning
$[0,0.1,0.2,0.4,0.8,1.5,3.0,100]\,\si{GeV}$ ends in a wide catch bin so that the
$\Eavail$-marginal captures the full CC-inclusive recoil tail; the global
OmniFold completeness is $c=1.0000$, as in 2D.

\subsection{The available-energy spectrum}
\label{sec:3d-spectrum}
Figure~\ref{fig:eavail} shows the unfolded $d\sigma/d\Eavail$ (the 1D projection
of the 3D result). It falls monotonically from
\SI{2.19e-38}{cm^2/GeV/nucleon} at the low-recoil peak to
\SI{6.8e-41}{} in the catch bin --- the expected CC-inclusive shape, with most
of the rate at $\Eavail\lesssim\SI{0.8}{GeV}$.

\begin{figure}[H]
  \centering
  \includegraphics[width=0.66\textwidth]{eavail_spectrum}
  \caption{Unfolded available-energy spectrum $d\sigma/d\Eavail$, the 1D
  projection of the 3D measurement. The wide $[3,100]\,\si{GeV}$ catch bin
  (value annotated) is excluded from the axis.}
  \label{fig:eavail}
\end{figure}

\subsection{Validation against the 2D measurement}
\label{sec:3d-anchor}
With no published 3D reference, the validation anchor is that marginalising the
3D cross section over $\Eavail$ must reproduce the 2D result of
\S\ref{sec:results}. The marginal recovers the 2D \emph{normalization} closely:
total $\sigma$ agrees to $+\SI{0.95}{\percent}$ and the per-bin
marginal/2D ratio has median $1.0016$. The \emph{shape} agreement is looser ---
the marginal versus the published covariance gives $\chi^{2}/\mathrm{ndf}=4.98$,
above the frozen 2D values ($3.66$ default, $2.65$ LightGBM-CV), driven by a
$\sim\SI{4.4}{\percent}$ per-bin scatter that the third axis injects into the
$(\pT,\pPar)$ marginal (Fig.~\ref{fig:eavail-marg}).

\begin{figure}[H]
  \centering
  \includegraphics[width=0.7\textwidth]{eavail_marginal_vs_paper_pull_full}
  \caption{Per-bin pull of the $\Eavail$-marginal against the published
  covariance, on the $14\times16$ grid and its distribution.}
  \label{fig:eavail-marg}
\end{figure}

\subsection{Closure: the scatter is physical, not a method bias}
\label{sec:3d-closure}
A reweight closure isolates whether that $\SI{4.4}{\percent}$ scatter is method
bias or real information. We inject a known $+\SI{30}{\percent}$ Gaussian bump
in \emph{truth} $\Eavail$ (center \SI{0.3}{GeV}, width \SI{0.15}{GeV}) into the
pseudo-data and require the 3D unfold to recover it. It does: the recovered
step-2 mean weight is $1.048$ against the injected $1.0485$, and the residual of
the unfold against the reweighted-truth reference is, on the $\Eavail$-marginal,
median $0.9999$ with standard deviation \textbf{\num{0.0006}}
(\SI{0.06}{\percent}); the per-3D-bin residual standard deviation is
\SI{3.2}{\percent} (MC statistics in sparse high-recoil corners). The closure
marginal scatter is thus $\sim70\times$ tighter than the
$\SI{4.4}{\percent}$ data-versus-2D scatter. The method is therefore unbiased on
the new axis, and we read the elevated marginal $\chi^{2}$ as
\emph{genuine data\,$\leftrightarrow$\,MC structure that the available-energy
dimension exposes} --- information that is invisible to the 2D measurement ---
rather than a pipeline artifact (the total normalization matches to
$<\SI{1}{\percent}$).

\subsection{Generator comparison in 3D}
\label{sec:3d-models}
Because the published ancillary ships only binned 2D model histograms, a
quantitative comparison of the 3D result (or of arbitrary projections) requires
generating truth-level events ourselves. No detector simulation is needed --- the
unfolded result already lives at truth level --- so a generator prediction is
obtained by histogramming its truth $d^{3}\sigma$ in the same fiducial phase
space and the same truth $\Eavail$ definition. We compare four predictions: \bpn{Did you run these simulators?}
\begin{itemize}
  \item \textbf{GENIE 2.12.10 central value} --- generated with \texttt{gevgen}
        on a CH target with the MINERvA ME~FHC flux and the
        DefaultPlusValenciaMEC configuration~\cite{Andreopoulos:2009rq}.
  \item \textbf{MINERvA Tune v1} --- the analysis tune (base GENIE $+$ RPA,
        low-recoil 2p2h, and non-resonant-pion reweights), extracted directly
        from the event-loop MC truth (the \texttt{w\_truth} weights), which
        reproduces the shipped ancillary curve to \SI{0.01}{\percent} in
        normalization.
  \item \textbf{NuWro 21.09}~\cite{Golan:2012wx} --- an independent generator
        (C target), the same flux and phase space.
  \item \textbf{GiBUU 2019}~\cite{Buss:2011mx} --- a transport-theory generator
        (C target), built and run \emph{natively} on the compute system (no
        container), same flux and phase space.
\end{itemize}
NEUT~\cite{Hayato:2009zz} is not included.
Each prediction is normalized by its own flux-averaged total CC cross
section per nucleon ($\langle\sigma\rangle\approx\SI{3.6}{}$--$\SI{4.0e-38}{cm^2}$).

Figure~\ref{fig:3dmodels} overlays the four on the unfolded result, with the
combined-covariance systematic band of \S\ref{sec:3d-syst} on the data. The
flux-averaged totals in the measurement phase space order as
GiBUU ($\SI{2.22e-38}{}$) $<$ NuWro ($\SI{2.34e-38}{}$) $<$ GENIE CV
($\SI{2.52e-38}{}$) $<$ Tune v1 ($\SI{2.71e-38}{}$) $<$ data
($\SI{3.08e-38}{cm^2/nucleon}$): all four models under-predict the rate, by
\SIrange{12}{28}{\percent}. On $(\pT,\pPar)$ the four track the data shape
closely; the discriminating axis is $\Eavail$, where NuWro, Tune v1 and GiBUU
\emph{suppress} the low-recoil quasielastic peak below bare GENIE (RPA and
nuclear effects) --- yet the data sit \emph{above} all four in the
lowest-$\Eavail$ bins. This low-available-energy excess is invisible to the 2D
measurement and is a direct statement of the unbinned method's high-dimensional
reach.

\begin{figure}[H]
  \centering
  \includegraphics[width=0.99\textwidth]{generators_vs_unfolded_band}
  \caption{1D projections of the unfolded 3D cross section (black, with the
  combined-covariance systematic band) versus four generator predictions
  produced at truth level: GENIE 2.12 CV, MINERvA Tune v1, NuWro 21.09, and
  GiBUU 2019. On $\pT$ and $\pPar$ the models track the data shape; on $\Eavail$
  they split, and all under-predict the data at low available energy. The wide
  $\Eavail$ catch bin is omitted from the axis.}
  \label{fig:3dmodels}
\end{figure}

Quantitatively, against the combined 3D covariance (\S\ref{sec:3d-syst}) the
integrated $\Eavail$ rate (catch bin dropped) places the models at
$-\SI{7.2}{\percent}$ (GENIE CV, $1.3\sigma$), $-\SI{9.5}{\percent}$ (Tune v1,
$1.8\sigma$), $-\SI{15.3}{\percent}$ (NuWro, $2.9\sigma$) and
$-\SI{21.9}{\percent}$ (GiBUU, $4.1\sigma$) below the data. A full-3D
goodness-of-fit over all 1431 bins, evaluated with a truncated-spectral
$\chi^{2}$ on the rank-247 constrained subspace of the combined covariance (a
raw pseudo-inverse of the rank-deficient matrix is never taken), excludes all
four at $p\approx0$ but ranks them sharply: diagonal-only $\chi^{2}/\mathrm{ndf}$
is Tune~v1 \textbf{4.8}, GiBUU 32.4, GENIE~CV 34, NuWro 36; the full
correlated metric is Tune~v1 131, NuWro 1287, GENIE~CV 1512, with GiBUU worst of
all because \SI{23.5}{\percent} of its residual lies \emph{outside} the
constrained subspace (its 3D shape differs in directions the covariance barely
constrains) on top of the largest normalization offset. The truncation scan and
per-generator eigen-mode residuals are shown in Fig.~\ref{fig:3dfullcov}.

\begin{figure}[H]
  \centering
  \includegraphics[width=0.9\textwidth]{compare_3d_fullcov}
  \caption{Full-3D truncated-spectral goodness-of-fit: per-generator
  $\chi^{2}/\mathrm{ndf}$ versus eigenvalue-truncation tolerance, with the
  fraction of the residual captured by the constrained subspace. The ranking
  (Tune~v1 best, GiBUU worst) is stable across the truncation choice.}
  \label{fig:3dfullcov}
\end{figure}

\subsection{The 3D systematic-uncertainty budget}
\label{sec:3d-syst}
The band on Fig.~\ref{fig:3dmodels} and the $\chi^{2}$ above rest on a full 3D
analog of the 2D uncertainty budget of \S\ref{sec:systematics}. The
systematic universes are dumped on the three-axis
$(\pT,\pPar,\Eavail)$ grid (187 universes), the bootstrap
statistical and ML-seedscan covariances propagated, and the three summed into a
single combined
covariance $C_{\mathrm{syst}}+C_{\mathrm{stat}}+C_{\mathrm{ML}}$ over all 1431
reported bins. The result has $\sqrt{\operatorname{tr}}=\SI{5.72e-39}{}$, a
median per-bin relative uncertainty of \SI{10.4}{\percent}, and numerical rank
$247$ (of 1431); it is flux-dominated, with the same systematic band ordering as
the 2D measurement (Fig.~\ref{fig:3dbands}). This combined covariance is what
projects onto each axis for
the per-bin bands and onto the full space for the goodness-of-fit, and it is the
internally-correct reference for the ``ours-only'' $\chi^{2}$ given that our
bootstrap statistical block is genuinely smaller than the paper's (an OmniFold
efficiency, confirmed by the data/MC split; see App.~\ref{sec:rank}).

\begin{figure}[H]
  \centering
  \includegraphics[width=0.49\textwidth]{uq_universe_3d_band_eavail}
  \includegraphics[width=0.49\textwidth]{uq_universe_3d_band_pt}
  \caption{Grouped fractional systematic bands of the 3D combined covariance
  projected onto $\Eavail$ (left) and $\pT$ (right): flux-dominated, with the
  same band ordering as the 2D budget.}
  \label{fig:3dbands}
\end{figure}

\subsection{The low-available-energy excess recovers the known low-recoil 2p2h deficit}
\label{sec:3d-2p2h}
The low-$\Eavail$ excess is the most physically interesting feature the third
axis exposes, and the 3D framework lets us identify its origin directly by
generating the relevant physics ourselves rather than only reweighting. It sits
in a well-established MINERvA low-recoil program: the enhancement of the event
rate between the quasielastic and $\Delta$ peaks, with multi-nucleon final
states, was first isolated in the low-three-momentum-transfer
subsample~\cite{MINERvA:2016lowrecoil} and is the effect the empirical MINERvA tune was
built to absorb --- our measurement re-exposes it on the muon-kinematics
$\times\,\Eavail$ axes.

\emph{Final-state interactions do not explain it.} A \texttt{grwght1p} reweight
pass on the 2M central-value GENIE events varying the pion-FSI dials shows both
are sub-percent on $d\sigma/d\Eavail$: \texttt{FrInel\_pi} shifts the spectrum by
$\le\SI{0.74}{\percent}$ (peak in $0.10$--$\SI{0.20}{GeV}$) and \texttt{FrAbs\_pi}
by $\le\SI{0.82}{\percent}$ (peak in the lowest bin), with the integrated rate
moving only $\pm\SI{0.02}{}$--$\SI{0.03}{\percent}$ at $\pm1\sigma$. A
$<\SI{1}{\percent}$ knob cannot bridge a $7$--$\SI{15}{\percent}$ excess, so the
excess points to the initial-state/nuclear model, not rescattering.

\emph{The shape and size match 2p2h.} A subtlety drives the diagnosis: our base
GENIE central value was generated with the default event-generator list, which
contains \emph{no} meson-exchange-current (2p2h/MEC) events
($\texttt{mec}=0$ for all $1.48\times10^{6}$ CC events; the modes are QE
\SI{11}{\percent}, RES \SI{22}{\percent}, DIS \SI{66}{\percent}, COH
\SI{0.9}{\percent}). Any 2p2h is therefore \emph{added}, not reweighted.
Decomposing the central-value $d\sigma/d\Eavail$ by interaction mode (per-bin
mode count-fractions $\times$ the bin value) and overlaying the data localises
the excess to the quasielastic--$\Delta$ dip: the $[0.10,0.20)$ bin is
$+3.9\sigma$ and $[0.20,0.40)$ is $+2.3\sigma$, while the QE-dominated
$[0,0.10)$ already matches ($+0.8\sigma$) and $[0.40,1.50)$ match to
$\pm0.4\sigma$. \SI{57}{\percent} of the integrated deficit sits at
$\Eavail\le\SI{0.4}{GeV}$, exactly where 2p2h lives, and closing it requires a
2p2h component $\approx\SI{43}{\percent}$ of the QE rate
(\SI{62}{\percent} locally) --- the standard empirical/Valencia-MEC magnitude,
in contrast to the sub-percent FSI dials (Fig.~\ref{fig:modedecomp}).

\emph{Enabling Valencia 2p2h confirms it.} We regenerated the GENIE central
value with the empirical 2p2h turned on (\texttt{--event-generator-list
Default+CCMEC}, 2M events). Because the Nieves--Simo--Vacas MEC channel, though
present in the splines, is absent from the \texttt{gspl2root} total-CC graph, the
sample is normalized by the non-MEC CC count (anchoring QE$+$RES$+$DIS$+$COH to
the known total and letting MEC add on top by $1/(1-f_{\mathrm{MEC}})$, with
$f_{\mathrm{MEC}}=\SI{2.87}{\percent}$). The added 2p2h lands squarely in the dip
and moves GENIE toward the data: the $[0,0.10)$ bin goes from $+0.8\sigma$ to
$-0.04\sigma$ (onto the data), $[0.10,0.20)$ from $+3.9\sigma$ to $+2.85\sigma$,
and $[0.20,0.40)$ from $+2.3\sigma$ to $+1.0\sigma$ (Fig.~\ref{fig:mec}). It
fills \SI{46}{\percent} of the data$-$CV gap in the dip and \SI{27}{\percent} of
the integrated deficit (CV $-\SI{7.2}{\percent}\to-\SI{5.2}{\percent}$), and
improves the full-3D $\chi^{2}/\mathrm{ndf}$ from $1512$ to $1145$
($-\SI{24}{\percent}$). It remains well short of Tune v1 ($131$), because Tune v1
layers MINERvA's empirical low-recoil 2p2h \emph{enhancement} and RPA on top of
the stock Valencia MEC --- i.e.\ stock Valencia 2p2h is real but under-strength
for MINERvA low recoil, which is precisely the known MINERvA-tune story and is
consistent with NuWro and GiBUU (native, un-enhanced 2p2h) sitting low as well.
The one feature 2p2h does not touch is a separate $+2.2\sigma$ excess in the
high-$\Eavail$ $[1.50,3.00)$ bin, a DIS-tail modeling issue resolved on the
$(\Eavail,W)$ plane in \S\ref{sec:eavailw}.

This low-available-energy deficit is the same feature MINERvA's dedicated low-recoil
inclusive measurement reports: Ascencio \emph{et al.}~\cite{MINERvA:2022incl}, which
measures $d^{2}\sigma/(dq_{3}\,d\Eavail)$ for $q_{3}<\SI{1.2}{GeV}$ at the same
$\langle E_\nu\rangle\sim\SI{6}{GeV}$, finds GENIE and NuWro underpredicting the data in
the low-$\Eavail$ region absent an enhanced 2p2h component. Our excess appears in the
muon-kinematics extension (integrated over all $q_3$) rather than at fixed $q_3$, so the
comparison is of the $\Eavail$ \emph{shape} and direction rather than absolute
normalization; the consistency confirms that the OmniFold 3D extension recovers the
established MINERvA low-recoil/2p2h feature as a by-product of the muon-kinematics
unfolding. A quantitative
shape comparison awaits the public \cite{MINERvA:2022incl} data release.

\begin{figure}[H]
  \centering
  \begin{minipage}{0.49\textwidth}
    \centering
    \includegraphics[width=\textwidth]{mode_decomp_eavail}
    \caption{Central-value $d\sigma/d\Eavail$ decomposed by interaction mode and
    overlaid with the unfolded data. The data excess is localised to the
    quasielastic--$\Delta$ dip ($\Eavail\le\SI{0.4}{GeV}$), the shape and size of
    a missing 2p2h component.}
    \label{fig:modedecomp}
  \end{minipage}\hfill
  \begin{minipage}{0.49\textwidth}
    \centering
    \includegraphics[width=\textwidth]{compare_mec_eavail}
    \caption{Re-generating GENIE with the empirical Valencia 2p2h enabled (red,
    stacked on the no-2p2h CV) fills part of the low-$\Eavail$ dip --- about half
    the data$-$CV gap, the remainder being the MINERvA-tune enhancement.}
    \label{fig:mec}
  \end{minipage}
\end{figure}

\subsection{Fourth axis: three-momentum transfer $q_{3}$}
\label{sec:4d}
The same one-feature-per-observable property lets us add a fourth axis at no
methodological cost. We choose the three-momentum transfer $q_{3}=|\vec{q}\,|$
because $(\Eavail,q_{3})$ is the plane in which the quasielastic, 2p2h and
resonance processes separate, and because it makes the comparison to MINERvA's
dedicated low-recoil measurement (Ascencio \emph{et al.}~\cite{MINERvA:2022incl},
$d^{2}\sigma/dq_{3}\,d\Eavail$) bin-compatible. Truth $q_{3}$ is the generator
momentum transfer (\texttt{CVUniverse::Getq3True()}); reco $q_{3}$ is the
calorimetric estimator built from the muon $Q^{2}$ and the recoil system
(\texttt{CVUniverse::RecoQ3()}). The event loop evaluates both accessors over
all twelve playlists, exactly as for $\Eavail$. Unlike $\Eavail$,
$q_{3}$ is \emph{not} invariant under the lateral systematic shifts, so the
shifted-$q_{3}$ universes are dumped per universe rather than reused from the
central value. The driver feature stack becomes $(\pT,\pPar,\Eavail,q_{3})$; the
$q_{3}$ binning $[0,0.2,0.4,0.6,0.8,1.2,2.0,100]\,\si{GeV}$ is deliberately
coarse (seven bins) and ends in a catch bin, for
$14\times16\times7\times7=10\,976$ four-dimensional bins (4830 populated).

\paragraph{Marginal closure.} Marginalising the 4D result over $q_{3}$ reproduces
the 3D $\Eavail$ spectrum and its low-recoil excess
(Fig.~\ref{fig:q3excess}, right): the $\Eavail$-marginal data/GENIE-CV ratio is
$1.21$ in the lowest bin, dips to $1.01$ across the quasielastic--$\Delta$
region, and rises again to $1.16$--$1.23$ in the high-$\Eavail$ DIS tail --- the
same shape established in 3D (\S\ref{sec:3d-2p2h}), now recovered as a by-product
of the higher-dimensional unfold (integrated data/GENIE-CV $=1.13$). This is the
4D$\to$3D analog of the 3D$\to$2D Jacobian anchor of \S\ref{sec:3d-anchor}.

\paragraph{What $q_{3}$ adds.} The new content is the resolved $(\Eavail,q_{3})$
plane (Fig.~\ref{fig:q3excess}, left). Integrated over $\Eavail$ the
$q_{3}$-marginal discrepancy is mild ($1.03$--$1.19$, center panel), but the
joint map is structured: at fixed $\Eavail$ in the quasielastic--$\Delta$ dip the
data/GENIE-CV ratio \emph{rises} with $q_{3}$, the central value over-predicting
the low-$q_{3}$, moderate-$\Eavail$ cells (ratio down to $0.66$) while the data
excess concentrates at low $\Eavail$ (all $q_{3}$) and toward high $q_{3}$. This
is precisely the correlation a muon-only or fixed-$q_{3}$ measurement cannot
expose, and it is the natural input for the bin-compatible Ascencio comparison.
We show it here at central-value level (the base GENIE event-generator list
contains no 2p2h, as in \S\ref{sec:3d-2p2h}); per-cell significances use the
combined 4D covariance and are therefore withheld pending its corrected
production.

\paragraph{Bin-identical comparison to the published low-recoil measurement.}
The Ref.~\cite{MINERvA:2022incl} supplemental data (44-cell
$d^{2}\sigma/(d\Eavail\,dq_{3})$ with its full covariance, distributed inside
the public arXiv source package) enables a direct quantitative cross-check.
Merging both measurements onto their maximal common grid --- edges present in
\emph{both} binnings, with our $(\pT,\pPar)$ marginalisation restricted to
$\pPar<\SI{20}{GeV/c}$ to mirror the Ascencio muon gate
($\theta_\mu<20^{\circ}$, $1.5<p_\mu<\SI{20}{GeV}$) --- gives two super-cells,
$\Eavail<\SI{0.4}{GeV}$ in $q_{3}\in[0.4,0.6)$ and $[0.6,1.2)\,\si{GeV}$: the
low-recoil region where the 2p2h physics lives. Our central result sits
\SI{9.2}{\percent} and \SI{6.3}{\percent} above Ascencio.  The former pulls,
full-covariance $\chi^{2}$, and consistency claim used the now-quarantined 4D
covariance and are withheld until its corrected projection is available
(Fig.~\ref{fig:ascencio}). Caveats: the two share MINERvA
flux/detector systematics, treated here as independent; the $p_\mu$-vs-$\pPar$
gate differs by $\lesssim\SI{6}{\percent}$ kinematic smearing at the
\SI{20}{GeV} boundary. A finer 44-cell comparison would require re-unfolding
onto the Ascencio edges (trivial for the unbinned method) \emph{and} a
systematic sweep on that binning; given the agreement above it is left to
future work.

\begin{figure}[H]
  \centering
  \includegraphics[width=0.85\textwidth]{ascencio_fullcov_compare}
  \caption{Bin-identical comparison to Ascencio
  \emph{et al.}~\cite{MINERvA:2022incl} on the maximal common
  $(\Eavail,q_{3})$ grid.  The central-value comparison remains current; the
  panel's uncertainty-derived metrics use the quarantined historical 4D
  covariance and must not be quoted pending the corrected marginalisation.}
  \label{fig:ascencio}
\end{figure}

\begin{figure}[H]
  \centering
  \includegraphics[width=0.99\textwidth]{q3_excess_projection}
  \caption{Data/GENIE-CV ratio of the 4D unfold projected onto the
  $(\Eavail,q_{3})$ plane (left; blank cells are kinematically empty), with the
  $q_{3}$ (center) and $\Eavail$ (right) marginals. The $\Eavail$ marginal
  reproduces the 3D low-recoil excess; the joint plane shows the data\,$-$\,MC
  discrepancy to carry genuine $q_{3}$ structure (the ratio rising with $q_{3}$
  across the quasielastic--$\Delta$ dip). GENIE central value, no 2p2h in the base
  event-generator list.}
  \label{fig:q3excess}
\end{figure}

\paragraph{Systematic covariance under remediation.} The 4D central result and
its lower-dimensional anchors remain valid, but the previous 187-universe
budget and unified-throw adoption are withdrawn. The replacement uses actual
minus/plus endpoints, MAT mean-centered $1/N$ blocks, one fixed estimator seed,
mean-centered joint throws with a separately reported mean shift, and no scalar
jitter subtraction. Statistical and ML replicas are also being regenerated.
Because dump-all lateral branches are CV-support-limited, a targeted
selection-complete active-universe campaign is bounding migration effects; full
five-band coverage remains the publication gate. No current 4D uncertainty
magnitude or unified/block ratio is asserted until the replacement artifact is
committed.
